\section{Objectives vs.\ Outcomes}
\label{sec:objectives-retrospective}

We initially designed MERIT with the intention to pivot technical recruitment towards demonstrated, multi-source evidence in an attempt
to combat reliance on a singular self-reported data source that is often subject to parsing errors and embellishment. The key contributions
to meet our objectives were extensible metrics, data fusion and glass-box explainability to combat the core limitations of the current
state of ATS systems. We will now cover the strengths and weaknesses of our approach to each of the three pillars.

\subsection{Pillar 1: Extensible evaluation}

\textbf{Worked:} Recruiter-configurable weights, semantic skill matching, and JD-driven metric priorities that align relatively
well with expert consensus. Linear single-column CVs parse well enough for downstream scoring.

\textbf{Failed:} Overall CV skill $F_1$ is 49.4\%, driven by a limited dictionary scoped only to programming languages and web frameworks from Devicon.
Multi-column layouts also fail when it comes to name detection, falling from 53.3\% to 13.3\%, and the JD skill $F_1$ is 40.81\%.
Extensibility is real, but it is often configured on mis-read documents, confirming ResumeBench's documented parsing 
fragility~\cite{emnlp2025resumebench} and Gugnani and Misra's~\cite{gugnani2020} concern that fixed, corpus-bound 
scoring limits role adaptability.

\subsection{Pillar 2: Multi-source fusion}

\textbf{Worked:} Study~01B helps to recover hidden talent (Jordan Smith, rank~\#10 $\rightarrow$~\#2). Study~07 gives the best observed expert
Spearman alignment from MERIT in one cohort ($\rho = 0.733$, $n=10$, exploratory). Study~01A shows large cardinal shifts when evidence validates
or contradicts claims ($+53.8\%$ / $-12.6\%$), partially validating Montandon et al.'s~\cite{montandon2021_roles} supplementary-GitHub hypothesis.

\textbf{Failed:} Carl Corp is penalised for empty GitHub despite significant tenure, reproducing Montandon et al.'s~\cite{montandon2021_roles} 
limitation when public repositories are absent.
Study~10 shows that MERIT fails when presented with
batch-wide contradictory evidence. Studies~08--09 show that no engine wins every hiring context (Section~\ref{subsec:spearman-synthesis}).
Fusion is strongly dependent on the presence of strong, reliable evidence that is accessible.

\subsection{Pillar 3: Glass-box explainability}

\textbf{Worked:} Shapley axioms hold where applicable (Study~11), negative attributions make penalties legible, and Study~14 shows recruiters
overturning bad shortlists upon evidence of fraud or stuffing in a formative pilot ($N=5$, mean SUS 81.5), addressing the transparency 
deficit~\cite{bogen2018,raghavan2020} directionally through contestable recruiter reports.

\textbf{Failed:} Explainability scales poorly, with an $\approx 5.6\times$ slowdown due to the introduction of Shapley values for each candidate.
With more data sources, the cost of explainability scales as $O(2^n)$, making it impractical for a large number of sources. Blind Mode still
intentionally exposes engine scores and tier labels, which could be considered a form of bias. Transparent breakdowns can instead anchor
automation bias rather than eliminate it (Section~\ref{subsec:construct-validity}).
We do not achieve Lo et al.'s~\cite{lo2025aihiring} scale of confirmatory human agreement, Study~14 remains a formative pilot only.
